\begin{abstract}
Contemporary artificial intelligence systems are predominantly \emph{epistemological mirrors}: they are trained on human-curated data, optimize human-defined objectives, and inherit human cognitive biases. This paper introduces the \textsc{Apeiron} Framework, a research programme aimed at constructing a fundamentally different kind of AI---one designed to autonomously generate knowledge from elementary observables without presupposing human categories, temporalities, or causal structures.

The framework is structured as a seventeen-layer architecture, organized into four conceptual clusters: Substrate \& Emergence (Layers 1--4), Adaptive Intelligence (Layers 5--7), Ontological Fluidity (Layers 8--13), and World-Making \& Governance (Layers 14--17). In this inaugural paper, we provide a complete formal specification and a reference implementation for Cluster I, establishing the irreducible observables (Layer~1), relational hypergraph (Layer~2), functional clusters (Layer~3), and endogenous temporal dynamics (Layer~4) that constitute the foundational substrate. Clusters II--IV are presented as detailed architectural specifications that define the trajectory for future theoretical and empirical work.

The contribution of this paper is threefold: (i) a rigorous, multi-axial definition of \emph{atomicity} grounded in Boolean algebra, measure theory, category theory, and information theory; (ii) a layered ontology with explicit translation functors that operationalize weak emergence and functional autonomy, validated through a non-trivial experiment on the MNIST dataset where the system spontaneously discovers functional clusters from raw pixels without labels; and (iii) a suite of falsifiable predictions that chart the course for progressive validation of the entire architecture. We do not claim that the seventeen-layer system is fully implemented or empirically validated; rather, we offer a coherent research programme and an open invitation to the scientific community to critique, refine, and collaboratively develop a genuinely non-anthropocentric artificial intelligence.
\end{abstract}